\documentclass[11pt]{article}
\usepackage[margin=.85in]{geometry}
\usepackage{amsmath,amssymb,hyperref}
\hypersetup{colorlinks=true,urlcolor=blue}
\newcommand{\Z}{\mathbb Z}
\newcommand{\F}{\mathbb F}
\newcommand{\Q}{\mathbb Q}
\newcommand{\up}{\operatorname{unit}_p}
\title{\vspace{-1.8em}Sharp stabilization of the prime-to-\(p\) part\\ of spanning-tree counts}
\author{Research supplement}
\date{26 November 2026}
\begin{document}
\maketitle
\begin{abstract}
Let a connected regular graph tower have compact \(p\)-adic analytic
pro-\(p\) deck group \(\Gamma\), with levels defined by the subgroups
\(W^{p^n}\) of an open normal uniform subgroup. Strip all factors of
\(p\) from the spanning-tree count at level \(n\) to obtain \(U_n\).
Then \(U_n\equiv U_{n-1}\pmod{p^n}\), with \(2^{n+1}\) replacing \(2^n\)
for \(p=2,\ n\ge2\). These moduli are sharp at every level, even for
simple graph towers. Independent exact checks include nonabelian cases.
\end{abstract}

\section{Finite-step reduction}
For a nonzero rational \(z\), put \(\up(z)=z/p^{v_p(z)}\).
Let \(G=\Gamma/W^{p^n}\) and \(K=W^{p^{n-1}}/W^{p^n}\).
The cover \(X_n\to X_{n-1}\) arises from quotienting by \(K\).
The kernel \(K\) is elementary abelian and is central in \(W/W^{p^n}\).
Every element of \(K\) is the \(p^{n-1}\)-st power of an element of
\(W/W^{p^n}\), a standard property of uniform pro-\(p\) groups.
These kernel properties alone suffice for the proof: for each nontrivial
character of \(K\), it is enough to have one such power it detects,
with the root centralizing \(K\) and of order dividing \(p^n\).

Choose the splitting field \(L=\Q(\zeta_{p^M})\) where \(p^M\) is
a multiple of the exponent of \(G\) and \(M\ge n\). Let \(A\) be its
Galois group over \(\Q\), identified also with the local Galois group
at \(p\). Decompose the lifted graph Laplacian into its blocks, using
the regular representation of \(G\) on each fiber over the base.
For each irreducible representation \(\rho\), the block \(D_\rho\)
appears \(d_\rho=\dim\rho\) times. It is invertible for nontrivial
\(\rho\) because connectivity makes the constants the whole kernel.
The matrix-tree theorem cancels the factors from representations
trivial on \(K\) and gives
\[
 \frac{\kappa_n}{\kappa_{n-1}}=
 \frac{1}{|K|}\prod_{\rho:\,K\not\subseteq\ker\rho}(\det D_\rho)^{d_\rho}.
\]
All blocks are Hermitian in unitary bases; their determinants here
are real and positive. The Galois action on these determinants agrees
with its action on the representations, since the voltage Laplacian
has rational group-ring coefficients.
\newpage
\section{The stabilizer and the orbit norm}
Fix an irreducible \(\rho\) in this product, with stabilizer \(H\le A\)
of its character and dimension \(d_\rho=p^k\).
Restrict it to \(K\). Its character eigenspaces are permuted transitively
by \(G\), by irreducibility and normality, and only nontrivial
characters of \(K\) occur.

Write \(\bar H_n\) for the image of \(H\) modulo \(p^n\).
Every \(\sigma\in H\) acts trivially on the \(p\)-th roots of unity:
if \(\sigma\) sends a constituent \(\chi\) of \(\rho|_K\) to \(\chi^a\),
then \(\chi^a=g\chi\) for some \(g\in G\).
Iterating gives \(g^j\chi=\chi^{a^j}\); the order of \(g\) is
a power of \(p\), whereas \(a\in\F_p^\times\) has order dividing \(p-1\).
This forces \(a=1\). (For \(p=2\) this is automatic.)

Choose \(w\in W/W^{p^n}\) with \(\chi(w^{p^{n-1}})\ne1\).
It preserves the \(\chi\)-space and all its eigenvalues there are
primitive \(p^n\)-th roots. The multiset is \(H\)-invariant:
Galois conjugation preserves \(\chi\) and the isomorphism type of
\(\rho\), hence preserves the restricted multiset for \(w\).
The action of \(\bar H_n\) on the primitive roots is free. So
\(|\bar H_n|\) divides each such eigenspace dimension (choosing \(w\)
for that constituent), and divides \(d_\rho\). Write
\(|\bar H_n|=p^s\), so \(s\le k\).

Let \(\alpha=\det D_\rho\).
For odd \(p\), the orbit contribution is \(N_{L^H/\Q}(\alpha)^{d_\rho}\).
We use cyclotomic local reciprocity: for \(B\le A\), the unit part \(u\)
of a norm from \(L^B\otimes\Q_p\) belongs to \(B\) modulo \(p^M\).
Indeed the cyclotomic reciprocity map sends \(p\) to the identity
and units to their residue (or its inverse, with the opposite
convention), and norms from the fixed field map to \(B\).
Thus \(u\bmod p^n\in\bar H_n\subseteq 1+p\Z/p^n\Z\).
This cyclic group has a unique subgroup of order \(p^s\), so
\(u\equiv1\bmod p^{n-s}\); raising \(u\) to \(p^k\) gives
\(1\bmod p^{n-s+k}\), in particular \(1\bmod p^n\).

For \(p=2,\ n\ge2\), realness puts \(\alpha\) in \(L^J\), where
\(J=\langle H,-1\rangle\). Now the orbit contribution is
\(N_{L^J/\Q}(\alpha)^{2^k[J:H]}\).
Normalize this norm to its unit part \(u\), and write \(u=\pm v\) with
\(v\in1+4\Z_2\). The image of \(J\) contains \(-1\), so \(v\bmod2^n\)
is in its principal-unit subgroup, whose order is half the order of
the image. If \(J=H\), this subgroup has order \(2^{s-1}\).
Then \(k\ge s\ge1\), so the sign vanishes under the exponent and
\[
 v^{2^k}\equiv1\pmod{2^{\,n+1-s+k}}.
\]
If \(J\ne H\), the subgroup has order at most \(2^s\).
The exponent \(2^{k+1}\) removes the sign and gives the same
precision. These follow from the cyclic structure of
\((1+4\Z_2)/(1+2^n\Z_2)\) and the gain of one bit per squaring on
\(1+4\Z_2\).
At \(p=2,\ n=1\) the asserted congruence just says the units are odd.
Multiplying the orbit contributions proves the claimed congruences
for \(U_n/U_{n-1}\), hence for \(U_n-U_{n-1}\).

Consequently \(U_n\) converges in \(\Z_p^\times\).
The residue at level \(N\ge1\) gives the limit modulo \(p^{N+1}\)
at odd \(p\), and modulo \(2^{N+2}\) at \(p=2\).
\newpage
\section{Sharpness at every level}
Use the additive \(\Z_p\)-tower with \((r-1)^2\) voltage loops of step
\(1\) and \(r\) loops of step \(2\). Its nontrivial Laplacian factor,
with \(z\) a character value, is
\[
 (2-z-z^{-1})(r+z)(r+z^{-1}).
\]
For primitive roots of order \(p^n\) with even totient, the product
over the layer is \(p^2\Phi_{p^n}(-r)^2\). After the matrix-tree
division by \(p\) between successive levels, the unit ratio is
\(\up(\Phi_{p^n}(-r))^2\). The same square formula holds at order
2 by direct evaluation.

For odd \(p\), take \(r=p-1\) except \(r=5\) at \(p=3\).
Write \((-r)^{p^{n-1}}=1+x\); the lifting-the-exponent formula gives
\(v_p(x)=n\). The geometric sum yields
\[
 \frac{\Phi_{p^n}(-r)}p
   =1+\frac{p-1}{2}x+
        \sum_{j=2}^{p-1}\frac1p\binom{p}{j+1}x^j .
\]
Every term in the displayed sum has valuation at least \(n+1\):
the last has valuation \((p-1)n-1\), and the other terms have
valuation at least \(2n\). The one exception to this estimate is
\(p=3,n=1\), where \(\Phi_3(-5)/3=7\) and \(v_3(7^2-1)=1\).
In all cases \(\Phi_{p^n}(-r)/p\) is a unit and squaring gives
a unit ratio different from 1 by exactly valuation \(n\).

For \(p=2,\ r=3,\ n\ge2\), the cyclotomic unit is
\((3^{2^{n-1}}+1)/2\). Its difference from 1 has valuation \(n\)
and its sum with 1 has valuation 1. Its square differs from 1
by valuation \(n+1\).
For \(p=2,\ n=1\), three step-1 loops give unit ratio 3, sharp
modulo 2. Later levels are more divisible, so these examples also
show the stated error bounds for the limit are optimal.

To obtain \emph{simple} graphs at all levels, subdivide each edge
into a path of length \(p^2\), with separate internal vertices for
each edge, doing so equivariantly. This includes base loops, for which
the replacement is a cycle of length at least four. Replacing every
edge by length \(\ell\) multiplies the tree count by
\(\ell^{|E|-|V|+1}\): edges in a spanning tree of the original graph
use a full path, and every other path has exactly one of its \(\ell\)
edges deleted. Thus the subdivision does not change the unit parts.

\section*{Sources, comparisons, and checks}
The abelian context includes
\href{https://arxiv.org/abs/2210.06182}{Ueki--Yoshizaki on cyclic resultants},
\href{https://arxiv.org/abs/2503.06194}{Yoshizaki on iterated multivariable resultants},
and \href{https://imath.kyiv.ua/~topology/conf/agma2025/abstracts/texts/kobayashi/kobayashi.pdf}{Kobayashi's graph class-number program},
including the Fibonacci tower (the same step pattern at \(r=2\)).
This continues the \href{https://chatgpt.com/c/6b07d086-1ef0-83e8-9354-9bbcf6afcaf3?messageId=88092bdb-6383-595f-a45a-d402786a8c09}{earlier cyclotomic transfer},
but uses Galois stabilizers and norm groups instead of individual
cyclotomic valuations. In particular, the representation argument
replaces the scalar characters of the abelian resultant setting. No hypothesis about zeta zeros is needed.
The proof is not Lean-formalized; publication priority is not claimed.
The Python replay computes independent exact integer Laplacian
cofactors. At \(n=2\) a nonabelian 81-vertex example gives
\((U_2,U_1)\equiv(1,10)\bmod27\); a nonabelian 64-vertex example
gives \((13,5)\bmod16\). The decks and steps are specified in the code.
\end{document}
